% Generated draft: requires mathematical and source review.
\documentclass[11pt]{article}
\usepackage[margin=25mm]{geometry}
\usepackage{fontspec}
\setmainfont{Latin Modern Roman}
\usepackage{amsmath,amssymb,mathtools,unicode-math}
\setmathfont{Latin Modern Math}
\usepackage{enumitem,needspace,xurl,hyperref}
\setlength{\parindent}{0pt}
\setlength{\parskip}{5pt}
\setlength{\emergencystretch}{3em}
\newcommand{\N}{\mathbb{N}}
\newcommand{\Z}{\mathbb{Z}}
\newcommand{\Q}{\mathbb{Q}}
\newcommand{\R}{\mathbb{R}}
\newcommand{\C}{\mathbb{C}}
\newcommand{\F}{\mathbb{F}}
\newcommand{\A}{\mathbb{A}}
\newcommand{\PP}{\mathbb{P}}
\newcommand{\E}{\mathbb{E}}
\newcommand{\eps}{\varepsilon}
\newcommand{\dd}{\,\mathrm{d}}
\begin{document}
\section*{UnsolvedMath ID: 20000390}
\textbf{Source problem code:} catalogue entry
\subsection{Definitions and mathematical statement}
Fix a prime power $q$, a polynomial $h\in\F_q[x]$ and an integer $k\geq3$. The AIM question asks for asymptotics, with $q,h,k$ fixed and $n\to\infty$, for the shifted divisor correlation $\sum_{\deg f=n}d_k(f)d_k(f+h)$. The original does not specify monicity or unit conventions for $d_k$ and allows $h=0$ in its wording.

In the supplied report's monic normalization, $d_k(f)$ counts ordered factorizations of a monic $f$ into $k$ monic factors, including factors $1$. For fixed nonzero $h$ and $n>\deg h$, set
\[
 S_k(n;h)=\sum_{\substack{f\text{ monic}\\\deg f=n}}d_k(f)d_k(f+h),\qquad r=k-1.
\]
For a monic polynomial $L$, put $|L|=q^{\deg L}$. The exactly solvable small-lcm sector is
\[
 T_{r,h}(n)=\sum_{\substack{A,B\text{ monic}\\\deg\operatorname{lcm}(A,B)\leq n\\\gcd(A,B)\mid h}}
 \frac{d_r(A)d_r(B)}{|\operatorname{lcm}(A,B)|}.
\]
The compatibility condition is that the congruences $f\equiv0\pmod A$ and $f\equiv-h\pmod B$ be simultaneously soluble. The report claims
\[
 S_k(n;h)\geq q^n T_{r,h}(n),\qquad
 T_{r,h}(n)=\frac{C_{q,r}(h)}{(2r)!}n^{2r}+O_{q,k,h}(n^{2r-1}),
\]
with a convergent positive shift-sensitive Euler product. In this CRT normalization its local factors can be written using monic irreducibles $P$, $Q_P=q^{\deg P}$ and $v_P(h)$ as
\[
 C_{q,r}(h)=\prod_P\left[(1-Q_P^{-1})^{2r}
 \sum_{\substack{a,b\geq0\\\min(a,b)\leq v_P(h)}}
 \binom{a+r-1}{r-1}\binom{b+r-1}{r-1}Q_P^{-\max(a,b)}\right].
\]
This is the product associated to the displayed sector; the report's detailed convergence and coefficient analysis is not supplied.

Its consequence is the positive lower bound
\[
 \liminf_{n\to\infty}\frac{S_k(n;h)}{q^n n^{2k-2}}
 \geq\frac{C_{q,k-1}(h)}{(2k-2)!}>0.
\]
This only bounds part of the sum from below. Contributions with lcm degree exceeding $n$ are not controlled here, so no matching upper bound or full asymptotic is asserted. The zero-shift moment and a sum over all leading coefficients require separate conventions and analysis.
\subsection{Short English statement}
Find fixed-$q$ asymptotics for higher divisor correlations at a fixed shift. The report handles monic polynomials and nonzero shifts, computing the complete small-lcm CRT sector and a positive lower bound on the expected scale. It does not determine the entire correlation or cover the zero shift.
\subsection{Sources}
\begin{itemize}
\item Local registry: \url{data/unsolvedmath/problems.json}, permanent ID $20000390$, source code \texttt{catalogue entry}; original question and supplied research summary.
\item AIM workshop ``Arithmetic statistics over finite fields and function fields'', Question 12, attributed to Julio Cesar Bueno de Andrade: \url{https://aimath.org/pastworkshops/arithstatffieldproblems.pdf}. Canonical location \texttt{aim-algebraic-number-theory-notes.json}, \texttt{notes[41]}; original remarks and literature field empty. No PDF was opened here.
\item Supplied machine-generated research attempt 1: partial result; problem partially solved at the run and in the registry; rigor metadata ``proved'', novelty confidence low. No full authoritative citation or proof of the sector Euler product/asymptotic is supplied. This non-peer-reviewed aid requires independent verification before citation.
\end{itemize}
\end{document}
